% IFN680 Project 5 technical report.
% Generated by tools/build_report.py. Every number below is printed by main_report.ipynb.
% Build:  tectonic project5_report.tex
\documentclass[10pt,a4paper,twocolumn]{article}

\usepackage[margin=1.5cm,top=0.9cm,bottom=0.8cm]{geometry}
\usepackage{lmodern}
\usepackage[T1]{fontenc}
\usepackage[utf8]{inputenc}
\usepackage{microtype}
\usepackage{graphicx}
\usepackage{booktabs}
\usepackage[font=small,labelfont=bf,skip=3pt]{caption}
\usepackage{titlesec}

\titlespacing*{\section}{0pt}{4pt}{0.5pt}
\titleformat{\section}{\normalfont\bfseries}{\thesection.}{0.5em}{}
\setlength{\parskip}{0pt}
\setlength{\parindent}{1em}
\setlength{\columnsep}{0.6cm}
\setlength{\textfloatsep}{3pt}
\setlength{\intextsep}{3pt}
\setlength{\floatsep}{3pt plus 1pt minus 1pt}
\setlength{\tabcolsep}{4pt}
\pagestyle{empty}

\renewcommand{\topfraction}{0.94}
\renewcommand{\bottomfraction}{0.6}
\renewcommand{\textfraction}{0.06}
\renewcommand{\floatpagefraction}{0.75}

\begin{document}

% The banner sits inside the \twocolumn optional argument. In the article class \twocolumn[...]
% claims the page's double column top region, so a figure* declared afterwards can never be
% placed on page 1 and is silently deferred to page 2.
\twocolumn[{%
\begin{center}
{\large\bfseries Project 5: Extending an Addition LLM to Subtraction and Reverse-Order Prediction}\\[2pt]
{\bfseries Group 4}
\end{center}
\vspace{-2pt}
Student 1: Karan Rooprai \hfill Student 1 ID: n12498122 \\
Student 2: Nhu Hieu Nguyen \hfill Student 2 ID: n12194778 \\
\vspace{3pt}

\begin{center}
\includegraphics[width=0.98\textwidth]{figures/figure_1_overview.pdf}
\end{center}
\vspace{-4pt}
\begin{quote}
\small\textbf{Figure 1.} \textbf{Left:} validation sequence accuracy per epoch for both
directions, at the full training set (solid) and at 80,000 examples
(dashed). \textbf{Centre and right:} error counts rather than accuracies, because every accuracy
here exceeds 0.98 and rates render the two models as identical bars. Forward makes
7 errors in 10000 examples, Reverse 44, and they
fall in different places.
\end{quote}
\vspace{2pt}
}]

The Week 7 workshop trains a small causal Transformer to add by next-token prediction, completing
a prompt such as \texttt{23+1=} with \texttt{24}. Task 1 extends it to subtraction with signed
results, Task 2 trains a second model that emits the answer digits right to left, and Task 3
compares them on one shared held-out set. The tokeniser, architecture, causal-masked objective
and training loop follow the workshop; only the changes described below differ.

\section{Task 1: handling subtraction}
Two changes were needed. First, the vocabulary gains a single \texttt{-} token with two roles: the
subtraction operator in a prompt (\texttt{50-73=}) and the sign of a negative result
(\texttt{-23}). Reusing one symbol keeps the vocabulary at fifteen tokens and lets the model infer
from position which role is meant, which is something self-attention can represent and a second
token would have obscured. Second, the generator samples an operation as well as two non-negative
operands of at most three digits; Python's \texttt{str} of a negative integer already produces the
signed target, so no special case is needed when $b > a$. Training data is a balanced mix, half
addition and half subtraction, with unique prompts and disjoint train, validation and test splits
drawn by slicing one shuffled list. Because borrows and signs make subtraction harder than
addition, more epochs and a validation set were used to confirm convergence.

\section{Task 2: reverse-order prediction}
Only the target is reversed; the prompt is untouched, so both models receive identical inputs and
differ solely in output ordering. The answer string is reversed literally, so \texttt{31} becomes
\texttt{13} and the model emits the units digit first, matching the direction in which carries and
borrows propagate by hand. The brief does not define how to reverse a negative result, so the same
literal rule is applied: \texttt{-123} becomes \texttt{321-}, placing the sign last. That is
consistent with right-to-left computation, since the sign of a difference is only determined once
the magnitude is. At evaluation the generated string is un-reversed before being parsed. Keeping
the tokeniser, architecture and optimiser identical makes prediction direction the only variable
in the comparison.

\section{Task 3: evaluation and comparative analysis}
\textbf{Set-up.} Both models are evaluated on one shared held-out set of 10000 examples,
balanced between addition and subtraction (subtraction fraction 0.49). The set is
regenerated from the random seed inside \texttt{main\_report.ipynb} and checked against the
shipped file, and its intersection with the training and validation prompts is confirmed to be
empty. Decoding is greedy, so the figures reproduce exactly on CPU or GPU.

\textbf{Operation robustness.} Both models clear the expected level comfortably: Forward reaches
0.9993 overall (95\% CI 0.9986 to 0.9997) and Reverse
0.9956 (0.9941 to 0.9967). Split by operation,
Forward reaches 1.0000 on addition over 5076 examples and
0.9986 on subtraction, and Reverse 0.9947 and
0.9965 respectively (Table~1).

\textbf{Digit-level performance.} Accuracy is reported per place value over the answers that
actually have that place, rather than over zero-padded answers. The distinction matters: only
2490 of 10000 answers reach the thousands column, so padding would score the
other 9045 as correct there by construction and drive every high position to 1.000.
Measured that way, Forward is exact at the thousands, hundreds and tens and scores 0.9993 at the units, while Reverse is exact at the thousands and scores 0.9996 at the hundreds, 0.9959 at the tens and 0.9991 at the units. Sign accuracy is 1.0000 for Forward and 1.0000 for Reverse, so the sign-last convention of Task 2 costs nothing.

\textbf{Direction effects.} The two models were scored on identical examples, so the comparison is
paired and McNemar's exact test applies. Forward is right where Reverse is wrong 44
times, Reverse right where Forward is wrong 7 times, and the two agree on
9949 of 10000 examples; $p = 1.212e-07$. The accuracy difference is therefore real rather than sampling noise, though it is small in absolute terms.
Trained on the
full set both directions converge, reaching 0.95 validation sequence accuracy at epoch
8 and 19 respectively. The reverse curve is the more striking of the two: it stays near a tenth of its final accuracy until epoch 18, then passes both the expected level and 0.95 within 1 epoch of doing so (Figure~1, left). A transition that abrupt is a barrier being crossed rather than capacity being filled, which is also why the reduced reverse run, given fewer updates over the same number of epochs, never crosses it. Training-set size is the second axis. At 80,000 examples the two orderings finish at 0.9945 and 0.0979 final validation sequence accuracy, against 0.9997 and 0.9949 on the full set (Figure~1, left). The two set sizes are not directly comparable in epochs, and where they can be compared in optimiser updates they agree: Forward passes 0.95 after 12,000 updates on the reduced set against 12,000 on the full one. Reverse does not reach 0.95 on the reduced set within the epochs run, which is what a shorter run looks like as much as a smaller one.

\textbf{Error patterns.} The residual errors are not scattered. Each direction concentrates them in a different region of the problem. The two sets have no case in common.

7 of Forward's 7 errors fall on answers of a single digit, a group
holding only 92 of the 10000 examples, and 7 of them are wrong by
exactly one unit. Predicting left to right requires committing to the answer's length before
emitting any digit, so a three-digit subtraction that collapses to a single digit is the case that
ordering makes hardest.

42 of Reverse's 44 errors fall on the 92
prompts whose first operand is 2 digits longer than the second, a rate of
0.4565 there against 2 errors across the other 9908
examples. Of those failures 38 are wrong at the tens digit,
the first place the model reaches after the shorter operand has run out of digits.
Emitting units first makes the units column trivial, because both operands end there, but it
forces the model to detect that one operand is exhausted and carry the remainder
alone. That is the alignment reverse ordering makes hardest, and it is the complement of Forward's weakness.

One split reads like a counterexample. Reverse is less accurate on additions that carry nowhere (0.9894 over 850 examples) than on additions that do carry (0.9957 over 4226), which inverts the expectation that carrying is the harder case. The operand width table resolves it: an addition that carries in no column tends to have a small second operand, so the no-carry group is enriched in exactly the width mismatch above. Carrying is not what this model finds difficult; keeping the columns aligned is.

\begin{table}[t]
\centering
\footnotesize
% Six columns and a spelled out interval need more room than the default gap leaves.
% tabcolsep is set inside the table environment, so it applies to this table only.
\setlength{\tabcolsep}{3pt}
\begin{tabular}{@{}llrrrr@{}}
\toprule
Mode & Split & $n$ & Errors & Accuracy & 95\% CI \\
\midrule
Forward & overall & 10000 & 7 & 0.9993 & 0.9986 to 0.9997 \\
Forward & addition & 5076 & 0 & 1.0000 & 0.9992 to 1.0000 \\
Forward & subtraction & 4924 & 7 & 0.9986 & 0.9971 to 0.9993 \\
Reverse & overall & 10000 & 44 & 0.9956 & 0.9941 to 0.9967 \\
Reverse & addition & 5076 & 27 & 0.9947 & 0.9923 to 0.9963 \\
Reverse & subtraction & 4924 & 17 & 0.9965 & 0.9945 to 0.9978 \\
\bottomrule
\end{tabular}
\caption{Accuracy with the count behind it and a Wilson interval. Both directions clear the
expected level by a wide margin, so the informative object is the small number of errors rather
than the rates.}
\end{table}

\begin{table}[b]
\centering
\footnotesize
\begin{tabular}{@{}lrrl@{}}
\toprule
Prompt & True & Predicted & Failure \\
\midrule
\multicolumn{4}{@{}l}{\emph{Forward: short answer, long operands}}\\
505-496 & 9 & 8 & single digit slip \\
487-479 & 8 & 7 & single digit slip \\
199-200 & -1 & -2 & single digit slip \\
907-899 & 8 & 7 & single digit slip \\
\multicolumn{4}{@{}l}{\emph{Reverse: operands of very different width}}\\
309-0 & 309 & 319 & single digit slip \\
195+2 & 197 & 187 & single digit slip \\
809+1 & 810 & 800 & single digit slip \\
730-8 & 722 & 733 & 2 digit slips \\
\bottomrule
\end{tabular}
\caption{Representative failures of each model, taken from the full error list printed by
\texttt{main\_report.ipynb}. The two sets have no case in common.}
\end{table}

\section{Limitations and recommendations}
Several limitations bound what these results support. Each model was trained once, from one
seed, so
the difference in convergence speed between the directions is a single observation rather than a
distribution; the accuracy comparison is on firmer ground because it is paired and tested.
Operands are capped at three digits, so nothing here shows whether either direction generalises to
longer arithmetic, and the error analysis suggests that is exactly where they would diverge.
Positions are encoded absolutely, following the workshop, which is the representation most likely
to be responsible for the alignment failure identified above. The data-efficiency ablation also held the number of epochs fixed rather than the number of optimiser updates, which is what made a difference in passes over the data look like a difference in how much data the task needs; a design that fixes updates would test data volume directly.

The recommendations follow from those. Train several seeds before treating the convergence gap as
a property of the ordering rather than of one optimisation trajectory. Extend the operand range
and re-measure, because the failure modes found here are both about structure, answer length for
Forward and operand width for Reverse, and both should grow with digits. Test a relative or
learned positional encoding against the absolute one, since the reverse model's errors cluster at
exactly the place where relative positions would carry the alignment information it appears to
lack. Finally, if one direction has to be chosen for this task, choose Forward: it is more accurate on this test set by a margin the paired test separates from noise, it trains at least as easily, and its failure mode is the rarer of the two.

\end{document}
